% GENERATED FILE. Edit canonical lessons, then run npm run generate:books.
% canonical-source-hash: fnv1a64:e6f30ff7369e10ac
% canonical-lessons: GU-C192-hisab, GU-C192-kul, GU-C192-rakam, GU-C192-bil, GU-C192-vedaphvu

\chapter{Adding up the bill}
\label{ch:gu-adding-up-the-bill}
\hlchaptermodality{192}

\begin{chapteropening}
\textbf{By the end of this chapter:} \emph{I can read the total on a bill and the amount to pay.}

Complete the last lesson of chapter 192: I can read the total on a bill and the amount to pay.
\end{chapteropening}

\section[\texorpdfstring{\gu{હિસાબ} (hisāb)}{હિસાબ (hisāb)}]{\gu{હિસાબ} (hisāb) — reckoning, accounts}

\label{lesson:GU-C192-hisab}



\begin{quote}
Before the new one: say the Gujarati for profit, then the Gujarati for loss.
\end{quote}

\subsection*{You'll want to know: \gu{હિસાબ}}
\textbf{\gu{હિસાબ}} — \emph{hisāb} — \textquotedblleft{}reckoning, accounts\textquotedblright{}.

\begin{quote}
\textbf{\gu{દુકાનદાર} \gu{સાંજે} \gu{હિસાબ} \gu{કરે} \gu{છે}.} — \emph{dukāndār sā̃je hisāb kare che} — \textquotedblleft{}The shopkeeper does the accounts in the evening.\textquotedblright{}
\end{quote}

A masculine noun.

\begin{cousinweb}
\textbf{\gu{હિસાબ}} (\emph{hisāb}) is a loan from Arabic \emph{ḥisāb}, \textquotedblleft{}counting, reckoning\textquotedblright{}, which reached Gujarati through Persian. Its Arabic root is \emph{ḥ-s-b}, \textquotedblleft{}to count\textquotedblright{} (Lane).
\end{cousinweb}

\begin{grammarlens}[title={What you've built}]
One more word for everyday life at A2.
\end{grammarlens}

\subsection*{Your turn}
\begin{itemize}
\raggedright
  \item \emph{Say it:} \emph{hisāb}
  \item \emph{Say it:} \emph{hisāb}, once more
  \item \emph{Say it:} \emph{dukāndār sā̃je hisāb kare che} — The shopkeeper does the accounts in the evening.
  \item \emph{Recall:} say the Gujarati for profit, then the Gujarati for loss, then say \emph{hisāb} again
\end{itemize}

\begin{tcolorbox}[breakable,colback=teal!4,colframe=teal!35!black,title={Before you move on}]
What does \gu{હિસાબ} mean? (\textquotedblleft{}Reckoning, accounts\textquotedblright{}.) Say it once more.
\end{tcolorbox}
\section[\texorpdfstring{\gu{કુલ} (kul)}{કુલ (kul)}]{\gu{કુલ} (kul) — total, in all}

\label{lesson:GU-C192-kul}



\begin{quote}
Before the new one: say the Gujarati for loss, then the Gujarati for reckoning.
\end{quote}

\subsection*{You'll want to know: \gu{કુલ}}
\textbf{\gu{કુલ}} — \emph{kul} — \textquotedblleft{}total, in all\textquotedblright{}.

\begin{quote}
\textbf{\gu{કુલ} \gu{સો} \gu{રૂપિયા} \gu{થયા}.} — \emph{kul so rūpiyā thayā} — \textquotedblleft{}That came to a hundred rupees in all.\textquotedblright{}
\end{quote}

\begin{cousinweb}
\textbf{\gu{કુલ}} (\emph{kul}) is a loan from Arabic \emph{kull}, \textquotedblleft{}all, the whole\textquotedblright{}, which reached Gujarati through Persian. Its Arabic root is \emph{k-l-l} (Lane).
\end{cousinweb}

\begin{grammarlens}[title={What you've built}]
One more word for everyday life at A2.
\end{grammarlens}

\subsection*{Your turn}
\begin{itemize}
\raggedright
  \item \emph{Say it:} \emph{kul}
  \item \emph{Say it:} \emph{kul}, once more
  \item \emph{Say it:} \emph{kul so rūpiyā thayā} — That came to a hundred rupees in all.
  \item \emph{Recall:} say the Gujarati for loss, then the Gujarati for reckoning, then say \emph{kul} again
\end{itemize}

\begin{tcolorbox}[breakable,colback=teal!4,colframe=teal!35!black,title={Before you move on}]
What does \gu{કુલ} mean? (\textquotedblleft{}Total, in all\textquotedblright{}.) Say it once more.
\end{tcolorbox}
\section[\texorpdfstring{\gu{રકમ} (rakam)}{રકમ (rakam)}]{\gu{રકમ} (rakam) — amount, sum}

\label{lesson:GU-C192-rakam}



\begin{quote}
Before the new one: say the Gujarati for reckoning, then the Gujarati for total.
\end{quote}

\subsection*{You'll want to know: \gu{રકમ}}
\textbf{\gu{રકમ}} — \emph{rakam} — \textquotedblleft{}amount, sum\textquotedblright{}.

\begin{quote}
\textbf{\gu{આ} \gu{રકમ} \gu{સાચી} \gu{છે}.} — \emph{ā rakam sācī che} — \textquotedblleft{}This amount is correct.\textquotedblright{}
\end{quote}

\begin{cousinweb}
\textbf{\gu{રકમ}} (\emph{rakam}) is a loan from Arabic \emph{raqam}, \textquotedblleft{}mark, written figure\textquotedblright{}, which reached Gujarati through Persian. Its Arabic root is \emph{r-q-m}, \textquotedblleft{}to write, to mark\textquotedblright{} (Lane).
\end{cousinweb}

\begin{grammarlens}[title={What you've built}]
One more word for everyday life at A2.
\end{grammarlens}

\subsection*{Your turn}
\begin{itemize}
\raggedright
  \item \emph{Say it:} \emph{rakam}
  \item \emph{Say it:} \emph{rakam}, once more
  \item \emph{Say it:} \emph{ā rakam sācī che} — This amount is correct.
  \item \emph{Recall:} say the Gujarati for reckoning, then the Gujarati for total, then say \emph{rakam} again
\end{itemize}

\begin{tcolorbox}[breakable,colback=teal!4,colframe=teal!35!black,title={Before you move on}]
What does \gu{રકમ} mean? (\textquotedblleft{}Amount, sum\textquotedblright{}.) Say it once more.
\end{tcolorbox}
\section[\texorpdfstring{\gu{બિલ} (bil)}{બિલ (bil)}]{\gu{બિલ} (bil) — bill}

\label{lesson:GU-C192-bil}



\begin{quote}
Before the new one: say the Gujarati for total, then the Gujarati for amount.
\end{quote}

\subsection*{You'll want to know: \gu{બિલ}}
\textbf{\gu{બિલ}} — \emph{bil} — \textquotedblleft{}bill\textquotedblright{}.

\begin{quote}
\textbf{\gu{હોટેલનું} \gu{બિલ} \gu{મોટું} \gu{હતું}.} — \emph{hoṭelnũ bil moṭũ hatũ} — \textquotedblleft{}The hotel bill was big.\textquotedblright{}
\end{quote}

\begin{cousinweb}
\textbf{\gu{બિલ}} (\emph{bil}) is a loan from English \emph{bill}.
\end{cousinweb}

\begin{grammarlens}[title={What you've built}]
One more word for everyday life at A2.
\end{grammarlens}

\subsection*{Your turn}
\begin{itemize}
\raggedright
  \item \emph{Say it:} \emph{bil}
  \item \emph{Say it:} \emph{bil}, once more
  \item \emph{Say it:} \emph{hoṭelnũ bil moṭũ hatũ} — The hotel bill was big.
  \item \emph{Recall:} say the Gujarati for total, then the Gujarati for amount, then say \emph{bil} again
\end{itemize}

\begin{tcolorbox}[breakable,colback=teal!4,colframe=teal!35!black,title={Before you move on}]
What does \gu{બિલ} mean? (\textquotedblleft{}Bill\textquotedblright{}.) Say it once more.
\end{tcolorbox}
\section[\texorpdfstring{\gu{વેડફવું} (veḍaphvũ)}{વેડફવું (veḍaphvũ)}]{\gu{વેડફવું} (veḍaphvũ) — to waste}

\label{lesson:GU-C192-vedaphvu}



\begin{quote}
Before the new one: say the Gujarati for amount, then the Gujarati for bill.
\end{quote}

\subsection*{You'll want to know: \gu{વેડફવું}}
\textbf{\gu{વેડફવું}} — \emph{veḍaphvũ} — \textquotedblleft{}to waste\textquotedblright{}.

\begin{quote}
\textbf{\gu{તે} \gu{પૈસા} \gu{વેડફે} \gu{છે}.} — \emph{te paisā veḍaphe che} — \textquotedblleft{}He wastes money.\textquotedblright{}
\end{quote}

\begin{cousinweb}
The origin of \textbf{\gu{વેડફવું}} (\emph{veḍaphvũ}) is not settled. CDIAL has no entry for it, and EtymDB lists no source.
\end{cousinweb}

\begin{grammarlens}[title={What you've built}]
One more word for everyday life at A2.
\end{grammarlens}

\subsection*{Your turn}
\begin{itemize}
\raggedright
  \item \emph{Say it:} \emph{veḍaphvũ}
  \item \emph{Say it:} \emph{veḍaphvũ}, once more
  \item \emph{Say it:} \emph{te paisā veḍaphe che} — He wastes money.
  \item \emph{Recall:} say the Gujarati for amount, then the Gujarati for bill, then say \emph{veḍaphvũ} again
\end{itemize}

\begin{tcolorbox}[breakable,colback=teal!4,colframe=teal!35!black,title={Before you move on}]
What does \gu{વેડફવું} mean? (\textquotedblleft{}To waste\textquotedblright{}.) Say it once more.
\end{tcolorbox}
